\clearpage
\fancyhead[L]{\small ZACH. TABLES D'ABERRATION.}
\hypertarget{native-551-lower}{}\label{native-551-lower}
\begin{center}\small Göttingische gelehrte Anzeigen. Stück 196. S. 1956..1960. 7. Dezember 1812.\\[3pt]\rule{.72\linewidth}{.3pt}\end{center}

\emph{Nouvelles tables d'aberration et de nutation pour quatorze cent quatre étoiles, avec une table générale d'aberration pour les planètes et les comètes, précédées d'une instruction qui renferme l'explication de l'usage de ces tables, suivies de plusieurs nouvelles tables destinées à faciliter les calculs astronomiques, par le Baron} \textsc{de Zach}. De l'imprimerie de Madame Mine et Comp. Marseille. 52 und 136 Seiten in Oktav.

Seit der Entdeckung der Aberration und Nutation haben Astronomen auf vielerlei Weise versucht, die lästige Berechnung dieser kleinen Ungleichheiten zu erleichtern, die sich für beobachtende Astronomen täglich wiederholt. Dafür gibt es verschiedene, mehr oder weniger geschickt eingerichtete allgemeine Tafeln. Für einzelne Sterne gibt es auch besondere Tafeln, aus denen man die Aberration und Nutation nach der jeweiligen Stellung der Sonne und des Mondknotens direkt ablesen kann. Die Tafeln \SourceLeaf{552}der letzteren Art von \textsc{Mezger}, \textsc{von Zach} und \textsc{Cagnoli} sind bei allen praktisch tätigen Astronomen im Gebrauch. Sie erfassen aber noch so wenige Sterne, dass man meistens auf die allgemeinen Tafeln zurückgreifen muss. Herr \textsc{von Zach} hat nun in der vorliegenden Schrift versucht, die Rechnung für eine beträchtlich größere Zahl von Sternen noch etwas mehr zu erleichtern, als es selbst mit den zweckmäßigsten allgemeinen Tafeln möglich ist. Eigentliche Spezialtafeln für so viele Sterne wären viel zu umfangreich und kostspielig, weshalb er darauf verzichtet hat. Der Gedanke ist sehr einfach und wurde für die Aberration schon früher von anderen Astronomen genutzt. Die Aberration in der Rektaszension eines jeden Sterns lässt sich leicht in die Form $m\sin(\odot+A)$ bringen. Dabei bezeichnet $m$ die größte positive Aberration des Sterns, $\odot$ die jeweilige Länge der Sonne und $A$ die Ergänzung zu $360^\circ$ derjenigen Sonnenlänge, bei der die Aberration verschwindet und vom negativen zum positiven Wert übergeht. Die Größen $m$ und $A$ hängen von der mittleren Stellung des Sterns ab. Da es bei ihnen nicht auf sehr große Genauigkeit ankommt, lassen sich die einmal bestimmten Werte eine ganze Reihe von Jahren unverändert verwenden. Genauso verhält es sich mit der Aberration in Deklination, der Nutation in Rektaszension und der Nutation in Deklination, wobei $m$ und $A$ natürlich für jedes dieser Elemente andere Werte haben. Herr \textsc{von Zach} gibt diese vier Werte von $A$ auf Minuten und die Logarithmen der zugehörigen $m$ auf vier Dezimalstellen genau an. Die Angaben für 1404 der wichtigsten Sterne aus allen in Europa sichtbaren Sternbildern beruhen auf ihren mittleren Positionen für 1800 und nehmen 94 Seiten ein. Damit beschränkt sich die ganze Arbeit darauf, vier Sinuslogarithmen und vier Zahlen anhand ihrer Logarithmen nachzuschlagen. Man schlägt also die logarithmischen Tafeln an vier Stellen nach, einmal mehr als bei den \emph{allgemeinen} Tafeln von \textsc{Gauss}, muss aber insgesamt weniger schreiben. Daher wird man \textsc{von Zach}s neue Tafeln für alle darin verzeichneten Sterne den allgemeinen gerne vorziehen. In der sonst sehr lehrreichen Einleitung von Herrn \textsc{von Zach} hat sich in die theoretische Herleitung eine kleine Unrichtigkeit eingeschlichen. Der Hilfsgröße $A$ wird eine falsche Bedeutung zugeschrieben. --- Nach diesen Tafeln für 1404 Sterne folgen die 36 sogenannten \textsc{Maskelyne}-Sterne noch einmal. Zuerst stehen ihre Rektaszensionen nach \textsc{Maskelyne}, zusammen mit $A$ und $\log m$ in Zeit für Aberration und Nutation. Danach folgen die Deklinationen derselben Sterne nach \textsc{Piazzi}, zusammen mit den beiden darauf bezogenen Werten von $A$ und $\log m$, alles für 1802. Hier bemerken wir einige an sich freilich unbedeutende Abweichungen von der ersten Tafel, die sich durch die Änderung der Epoche allein nicht erklären lassen und Rechenfehler zu sein scheinen. So ergibt sich zum Beispiel bei der Deklination von $\alpha$ Tauri \[ \begin{array}{r@{}l|c|c} \multicolumn{2}{c|}{} & \text{erste Tafel} & \text{zweite Tafel}\\ \text{Aberration}\; & \left\{\begin{array}{l}\log m\\ A\end{array}\right. & \begin{array}{c}0\,5806\\7^{\mathrm s}\;23^\circ\;26'\end{array} & \begin{array}{c}0.5772\\7^{\mathrm s}\;23^\circ\;15'\end{array}\\ \text{Nutation}\; & \left\{\begin{array}{l}\log m\\ A\end{array}\right. & \begin{array}{c}0.9673\\3^{\mathrm s}\;18^\circ\;18'\end{array} & \begin{array}{c}0.9682\\3^{\mathrm s}\;18^\circ\;6'\end{array} \end{array} \]

Eine ähnliche Tafel folgt für die Sterne $\alpha$ und $\beta$ im Kleinen Bären nach ihren Stellungen in den Jahren 1790, 1800, 1810 und 1820.

Obwohl die Größen $m$ und $A$ eigentlich nur für das Jahr gelten, für das sie berechnet sind, kann man sie ohne erheblichen Fehler für eine ziemlich lange Reihe von Jahren vor und nach jener Epoche beibehalten. Der Verfasser gibt ein paar besondere Tafeln an, um abzuschätzen, um wie viel sich Rektaszension und Deklination ändern dürfen, bevor der dadurch entstehende Fehler auf 1 Sekunde anwächst. Diese Tafel hilft zugleich zu beurteilen, inwieweit \SourceLeaf{553}man für einen Stern, der nicht unter den 1404 der großen Tafel steht, die Konstanten eines benachbarten Sterns aus ihr verwenden darf. Einige darauf folgende Tafeln für Aberration, Nutation, Präzession und die Korrektur des Mittags aus entsprechenden Sonnenhöhen übergehen wir hier stillschweigend als bekannt, da der Verfasser sie schon anderswo mitgeteilt hat. --- Mit Tafel XX lässt sich die Sonnenlänge als Aberrationsargument für den Tag eines Jahres finden, für das man keinen astronomischen Kalender zur Hand hat. Sie gibt an, um wie viel die Epoche der mittleren Sonnenlänge in jedem einzelnen Jahr von 1700 bis 1827 größer ist als im Jahr 1803. Mit ausreichender Genauigkeit für den angegebenen Zweck kann man diese Zahlen zugleich als Unterschiede der wahren Längen an Tagen desselben Datums auffassen. So lassen sich die Sonnenlängen für jedes dieser Jahre leicht aus denen eines beliebigen anderen Jahres ableiten, für das man gerade einen astronomischen Kalender zur Hand hat. Ebenso gibt Herr \textsc{von Zach} noch eine Tafel für die Länge des Mondknotens zu Beginn aller Jahre von 1700 bis 1827, damit man die Ephemeriden eines beliebigen dieser Jahre auf ein anderes übertragen kann. Zum Schluss folgt eine Tafel zur Umrechnung der Sternzeit in mittlere Sonnenzeit. Sie ist noch etwas bequemer eingerichtet als die bekannten kleinen Tafeln des Verfassers, die für denselben Zweck berechnet wurden und bei allen Astronomen im Gebrauch sind. Beigefügt ist auch eine kleine Tafel für den Teil der Nutation, der von der Sonnenlänge abhängt. Diesen Teil hätte man gleich mit Tafel XXII verbinden können, was die Rechnung noch ein wenig verkürzt hätte. Vermutlich hat der Verfasser das absichtlich unterlassen, damit in dieser Tafel die reinen mittleren Rektaszensionen der Sonne stehen. Aus demselben Grund scheint er auf die Hinzufügung einer negativen Konstante verzichtet zu haben, die sonst verwendet wird, um die Mondnutation stets positiv zu machen.

Aus dieser kurzen Inhaltsangabe sieht man, dass sich der verdiente Verfasser mit der Sammlung von Tafeln erneut den Dank aller praktisch tätigen Astronomen erworben hat. Ihnen muss jede Erleichterung ihrer mühsamen, täglich wiederkehrenden Rechnungen sehr willkommen sein.
